\documentclass[output=paper]{langscibook}
\ChapterDOI{10.5281/zenodo.22880236}
\title{\textit{O overe} -- The coconut palm}
\author{Joice Maion}
\abstract{}
\begin{document}
\maketitle\label{chap:5}
\begin{figure}[H]
    \includegraphics[width=.9\linewidth]{figures/5_1.jpg}
    \caption{\textit{O overe} – The coconut palm}
    \label{fig:5_1}
\end{figure}
\begin{figure}[H]
    \includegraphics[width=.4\linewidth]{figures/5_2.png}
    \caption{\textit{Amaa hum no overe} -- The parts of the coconut palm}
    \label{fig:5_2}
\end{figure} 
\begin{xltabular}{\textwidth}{lXX}    
    (\ref{ex:text01-1})&\textit{\textbf{O vahutate no overe}}&\textbf{The story of the coconut palm}\\
    (\ref{ex:text01-2})&\textit{Bea overe mahaka maa gina teo overe vai to sun nana,}&When the coconut has become dry high up on a standing coconut palm,\\
    (\ref{ex:text01-3})&\textit{ore paa kuu gunaha vo maa teo kasuana.}&it falls down to the ground.\\
    (\ref{ex:text01-4})&\textit{Be obete teo kasuana bona buaku a sivao,}&When it is lying on the ground for two months,\\
    (\ref{ex:text01-5})&\textit{a vana repaa sigooru komana overe.}&the germinating leaves are sprouting inside the coconut.\\
    (\ref{ex:text01-6})&\textit{Saka tara haa raara a vana,}&We cannot see the germinating leaves,\\
    (\ref{ex:text01-7})&\textit{ei kou na tei nana komana overe.}&because they are inside the coconut.\\
    (\ref{ex:text01-8})&\textit{Evehee ean na nata kaakoo nom,}&But you know what is happening,\\
    (\ref{ex:text01-9})&\textit{bean tara o ana be sigooru tavusu.}&when you see the root coming out and sprouting.\\
    (\ref{ex:text01-10})&\textit{Beo ana sigooru tavusu,}&When the sprouting root comes out,\\
    (\ref{ex:text01-11})&\textit{eove repaa taneo tea pura hopo vo kasuana.}&it starts growing into the ground.\\
    (\ref{ex:text01-12})&\textit{Ean paa tara be boko tavusu maa o supi.}&You see that the sprouting leaf has broken out (of the coconut husk).\\
\end{xltabular}
\begin{figure}[H]
    \captionsetup[subfigure]{oneside,margin={0cm,0cm}}
    \subfloat[]{\includegraphics[trim=100 0 112 0,clip,width=.333\linewidth]{figures/5_3.jpg}\label{fig:5_3}}\hfil
    \subfloat[]{\includegraphics[trim=220 0 172 0,clip,width=.333\linewidth]{figures/5_4.jpg}\label{fig:5_4}}\hfil
    \subfloat[]{\includegraphics[trim=0 20 0 0,clip,width=.333\linewidth]{figures/5_5.jpg}\label{fig:5_5}}
    \caption{(a) \textit{O overe o manoto} -- The ripe coconuts, (b) \textit{O ana} -- The root, (c) \textit{O supi} -- The sprouting leaf}
\end{figure}
\begin{figure}[H]
    \includegraphics[trim=0 15 0 19,clip,width=.37\linewidth]{figures/5_6.jpg}\hfil
    \includegraphics[trim=0 0 0 21,clip,width=.35\linewidth]{figures/5_7.jpg}\hfil
    \includegraphics[trim=0 0 0 0,clip,width=.28\linewidth]{figures/5_8.jpg}
    \caption{\textit{O supi} -- The sprouting leaf}
    \label{fig:5_6-8}
\end{figure}
\begin{xltabular}{\textwidth}{lXX}
    (\ref{ex:text01-13})&\textit{Beo supi pura hiava bata,}&When the sprouting leaf grows upwards,\\
    (\ref{ex:text01-14})&\textit{o paka repaa taneo bata maa tea taparata.}&the leaf starts opening.\\
    (\ref{ex:text01-15})&\textit{Be taparata vabeera bata o paka no overe,}&When the leaves of the coconut palm have fully opened,\\
    (\ref{ex:text01-16})&\textit{o ana nae mee na pura hopo vakikisi batana teo kasuana.}&its roots also start growing into the ground, thereby giving the coconut a firm hold.\\
    (\ref{ex:text01-17})&\textit{Amaa meha maa ana mee paa taneou tea pura tavusu}&Other roots have also started growing out (of the nut)\\
    (\ref{ex:text01-18})&\textit{bara pura hopo vo teo kasuana.}&and growing into the ground.\\
    (\ref{ex:text01-19})&\textit{Beo overe sun vakikisi, be sa antee haana tea takara,}&When the coconut palm stands firmly, when it cannot be uprooted,\\
    (\ref{ex:text01-20})&\textit{o paka overe repaa takihi.}&the coconut leaf splits open.\\
    (\ref{ex:text01-21})&\textit{Beo meho supi tavusu, o paka vai to tavusu momohu, repaa takihi,}&When another sprouting leaf comes out, the leaf that came out first, splits open,\\
    (\ref{ex:text01-22})&\textit{ore pura hiava komana o overe,}&so that the coconut palm grows up,\\
    (\ref{ex:text01-23})&\textit{be paku bata vo na teebona.}&when it is happening like this.\\
    (\ref{ex:text01-24})&\textit{Beo overe pura hiava, a vana vaa komana overe repaa kavara.}&When the coconut palm has grown up, the germinating leaves inside the coconut are gone.\\
    (\ref{ex:text01-25})&\textit{Beo supi pihae pura tavusu, o iposi na rakoto batana teo pihae.}&When the sprout of the coconut frond grows outside, the leaf sheath sticks on the coconut frond.\\    
\end{xltabular}
\begin{figure}[H]
    \includegraphics[trim=0 100 0 0,clip,width=.33\linewidth]{figures/5_9.jpg}\hfil
    \includegraphics[trim=0 0 0 100,clip,width=.33\linewidth]{figures/5_10.jpg}\hfil
    \includegraphics[trim=0 0 201 0,clip,width=.34\linewidth]{figures/5_11.jpg}
    \caption{\textit{O iposi} -- Leaf sheath}
    \label{fig:5_9-11}
\end{figure}
\begin{xltabular}{\textwidth}{lXX}
    (\ref{ex:text01-26})&\textit{Beo pihae taparata, o iposi repaa taparata kahi bata bono pihae.}&When the coconut frond opens, the leaf sheath unfolds and separates from the frond.\\
    (\ref{ex:text01-27})&\textit{Beo pihae beera, o iposi repaa taparata vakavara,}&When the coconut frond becomes big, the leaf sheath completely opens,\\
    (\ref{ex:text01-28})&\textit{ore paa hana gunaha koa.}&so that it just hangs down.\\
    (\ref{ex:text01-29})&\textit{Beo pihae mahaka, ore paa kuu,}&When the coconut frond becomes dry, it falls down,\\
    (\ref{ex:text01-30})&\textit{o iposi repaa kuu bata mi bona.}&and then the leaf sheath falls down together with it.\\
\end{xltabular}
\begin{figure}[H]
    \includegraphics[width=.34\linewidth]{figures/5_12.jpg}\hfil
    \includegraphics[width=.31\linewidth]{figures/5_13.jpg}\hfil
    \includegraphics[width=.34\linewidth]{figures/5_14.jpg}
    \caption{\textit{O rakopo bara nanisi} -- The flower sheath and the panicle}
    \label{fig:5_12-14}
\end{figure}
\begin{figure}[H]
    \includegraphics[width=.5\linewidth]{figures/5_15.jpg}
    \caption{\textit{A kukusu} -- The very young coconut}
    \label{fig:5_15}
\end{figure}
\begin{figure}[H]
    \includegraphics[trim=0 0 0 0,clip,width=.55\linewidth]{figures/5_16.jpg}\hfil
    \includegraphics[trim=60 0 60 0,clip,width=.45\linewidth]{figures/5_17.jpg}
    \caption{\textit{A magura} -- The drinkable coconut}
    \label{fig:5_16-17}
\end{figure}
\begin{figure}[H]
    \includegraphics[width=.5\linewidth]{figures/5_18.jpg}
    \caption{\textit{O overe o manoto} -- The ripe coconuts}
    \label{fig:5_18}
\end{figure}
\begin{xltabular}{\textwidth}{lXX}
    (\ref{ex:text01-31})&\textit{O overe be vua, a vua overe na topo ge hana batana teo nanisi.}&The coconut palm, when it bears nuts, the coconuts stick or hang on the panicle.\\
    (\ref{ex:text01-32})&\textit{A peha totoka ge buaku totoka overe na antee nana tea hana teo peho nanisi overe.}&Six or seven coconuts can hang on one panicle.\\
    (\ref{ex:text01-33})&\textit{Beo overe vua, o rakopo na rakopo batana toon no ibu overe.}&When the coconut palm bears nuts, the flower-sheath covers the cluster of coconuts.\\
    (\ref{ex:text01-34})&\textit{Bea kukusu overe beera bata, o rakopo repaa takapa kahi bata bono ibu overe.}&When the very young coconut becomes bigger, the flower sheath peels off the cluster of coconuts.\\
    (\ref{ex:text01-35})&\textit{A kukusu overe, a overe vai a rutaa, a kobusu.}&The very young coconut is a small coconut that is watery inside.\\
    (\ref{ex:text01-36})&\textit{Ahiki vakisi ta tabaan komanae.}&There is not yet any food inside.\\
    (\ref{ex:text01-37})&\textit{Be beera eori repaa dao bari bona magura.}&When it grows bigger, they call it \textit{magura} ‘ripening coconut’.\\
    (\ref{ex:text01-38})&\textit{A pereperee ge kasikasiree a kapa nae.}&Its cover is green or yellow.\\
    (\ref{ex:text01-39})&\textit{A tabaan paa tavusu vaevuru komanae, evehee sabin kikisi.}&The food has already developed inside, but it is not hard yet.\\
    (\ref{ex:text01-40})&\textit{A tabaani nae a nin mataa kurus.}&Its food has a very good taste\\
    (\ref{ex:text01-41})&\textit{A huun vaa tea magura na gaagaga riori.}&The liquid in the \textit{magura} is drinkable.\\
    (\ref{ex:text01-42})&\textit{A huun bona a nini mataa kurusu.}&This liquid has a very good taste.\\
    (\ref{ex:text01-43})&\textit{Na mosi riori, eori repaa gaga bari.}&They cut (the \textit{magura}), then they can drink it.\\
    (\ref{ex:text01-44})&\textit{Bean tara a overe vai a manoto, a kapa nae na uruurunata nana.}&When you look at a ripe coconut, its cover is dark brown.\\
    (\ref{ex:text01-45})&\textit{Enam repaa boha voen a overe a mahaka.}&Then we say that the coconut is dry.\\
    (\ref{ex:text01-46})&\textit{Bea overe manoto, a tabaani vaa komanae paa kikisiu.}&When the coconut is ripe, the food inside has become hard.\\
    (\ref{ex:text01-47})&\textit{Bea overe manoto maa gina teo overe,}&When the coconut has become ripe high up on the coconut palm,\\
    (\ref{ex:text01-48})&\textit{ore paa kuu gunaha vo maa teo kasuana.}&it falls down to the ground.\\
    (\ref{ex:text01-49})&\textit{A overe vai a manoto na gono riori, eori repaa hovu bari.}&They get the ripe coconut, and then they husk it.\\
\end{xltabular}
\begin{figure}[H]
    \captionsetup[subfigure]{oneside,margin={0cm,0cm}}
    \subfloat[]{\includegraphics[trim=0 100 0 0,clip,width=.45\linewidth]{figures/5_19.jpg}\label{fig:5_19}}
    \subfloat[]{\includegraphics[trim=0 10 0 10,clip,width=.45\linewidth]{figures/5_20.png}\label{fig:5_20}}
    \caption{(a) \textit{Hovu overe} -- Husking coconuts, (b) \textit{O pusu hiovu} -- Husking stick}
\end{figure}
\begin{figure}[H]
    \includegraphics[width=.6\linewidth]{figures/5_21.jpg}
    \caption{\textit{Booboha overe} -- Breaking the shell of coconuts}
    \label{fig:5_21}
\end{figure}
\begin{xltabular}{\textwidth}{lXX}
    (\ref{ex:text01-50})&\textit{\textbf{Hovu}}&\textbf{Husking coconuts}\\
    (\ref{ex:text01-51})&\textit{Beara mei rara tea hovu overe, eara na vaamatatopo varaba rara o pusu hiovu.}&Before we husk coconuts, we prepare the husker.\\
    (\ref{ex:text01-52})&\textit{O pusu hiovu na paku raara bene ta boon.}&We make the husker from a piece of mangrove wood.\\
    (\ref{ex:text01-53})&\textit{Eara na mosi maa rara e ta boon tea boon.}&We cut a piece of mangrove wood in the mangroves.\\
    (\ref{ex:text01-54})&\textit{A bebeahu nae toro antee bona peha gea buaku a ato.}&Its length should be one or two yards.\\
    (\ref{ex:text01-55})&\textit{Eara repaa karem vakakaniva a peha papana.}&Then we sharpen one end.\\
    (\ref{ex:text01-56})&\textit{O papana bona to kahi hovu kaakoo raara.}&This end we will use for husking the coconut.\\
    (\ref{ex:text01-57})&\textit{A papana to sa karem vakakaniva vira haa raara,}&The end that we did not sharpen,\\
    (\ref{ex:text01-58})&\textit{eve to pasi tahi vaatopo ni noman teo kasuana.}&that is the one that you will push into the ground.\\
    (\ref{ex:text01-59})&\textit{Be kavara, ean repaa taneo tea hovu.}&When it is finished, you start husking.\\
    (\ref{ex:text01-60})&\textit{Ean na vaahuu hopo ni nom a overe tea ta hiovu.}&You stick the coconut onto the husking stick.\\
    (\ref{ex:text01-61})&\textit{Ean paa menaa bata a overe, o kapa repaa takapa kahi bata bono meho kapa.}&Then you bend the coconut so that the cover (the husk) peels off the other cover (the shell).\\
    (\ref{ex:text01-62})&\textit{A kukan bau gea dorana a hovu, a overe repaa tahovu vakavara.}&Perhaps three or four husking movements, and then the coconut is completely husked.\\
    (\ref{ex:text01-63})&\textit{O kapa to hovu noman na dao raara bono pesupaava}&The cover that you remove by husking is called \textit{pesupaava} `fibrous husk’.\\
    (\ref{ex:text01-64})&\textit{Nomana a pusu hiovu na paku riori bona maa hum aeana vai a bebeahu.}&Today they make the huskers from long pieces of iron.\\
    (\ref{ex:text01-65})&\textit{Na karem vakakaniva riori bona maa soo aeana.}&They sharpen (them) with iron files.\\
    (\ref{ex:text01-66})&\textit{Eori repaa hovu overe ni bari.}&Then they use them for husking coconuts.\\
    (\ref{ex:text01-67})&\textit{Ean he paa hovu vuru o kapa vaa gina, mean paa hovu kahi o meho kapa.}&When you have removed the top cover, you remove the other cover (the shell).\\
    (\ref{ex:text01-68})&\textit{O kapa bona to booboha raara, ei kou o kikisi.}&This cover (the nutshell), we break, because (it) is hard.\\
    (\ref{ex:text01-69})&\textit{A ba suburaara na booboha overe ni roho bona vasu.}&Our grandparents broke coconuts (by hitting them) with a stone,\\
    (\ref{ex:text01-70})&\textit{a overe repaa taboha.}&then the coconut broke.\\
    (\ref{ex:text01-71})&\textit{Evehee nomana na booboha overe ni rori bono rapisi,}&But today they break coconuts (by hitting them) with a bush knife,\\
    (\ref{ex:text01-72})&\textit{ore paa taboha.}&then they break.\\
\end{xltabular}
\begin{figure}[H]
    \includegraphics[trim=80 10 0 0,clip,width=.45\linewidth]{figures/5_22.jpg}\hfil
    \includegraphics[trim=20 0 20 3,clip,width=.55\linewidth]{figures/5_23.jpg}
    \caption{\textit{O korivahi} -- The Thorny Oyster}
    \label{fig:5_22-23}
\end{figure}
\begin{xltabular}{\textwidth}{lXX}
    (\ref{ex:text01-73})&\textit{\textbf{Kuukuhi overe}}&\textbf{Grating coconuts}\\
    (\ref{ex:text01-74})&\textit{Murinae eori repaa kuukuhi bari mene rovi bona maa nahu iana, paapao gea tabae to rake riori tea rovi.}&After that they grate them for cooking pots of fish, vegetables or whatever they want to cook in coconut cream.\\
    (\ref{ex:text01-75})&\textit{A ba suburaara na kuukuhi overe ni roho bona kehaa.}&Our grandparents grated coconuts with a shell.\\
    (\ref{ex:text01-76})&\textit{A hena no kehaa bona na dao riori bono korivahi.}&The name of this shellfish is \textit{korivahi} `Thorny Oyster’.\\
\end{xltabular}
\begin{figure}[H]
    \includegraphics[trim=0 0 0 123,clip,width=.35\linewidth]{figures/5_24.jpg}\hfil
    \includegraphics[width=.65\linewidth]{figures/5_25.jpg}
    \caption{\textit{Kuukuhi overe} -- Grating coconuts}
    \label{fig:5_24-25}
\end{figure}
\begin{figure}[H]
    \includegraphics[width=.55\linewidth]{figures/5_26.jpg}
    \caption{\textit{A kuhi} -- The grater}
    \label{fig:5_26}
\end{figure}
\begin{xltabular}{\textwidth}{lXX}
    (\ref{ex:text01-77})&\textit{O korivahi na nibusu riori,}&They made a hole into the Thorny Oyster,\\
    (\ref{ex:text01-78})&\textit{eori repaa nibusu pete bene ta naono to paa kahu vamatatopo vuru ori.}&then they also made a hole into a piece of wood that they had already carved.\\
    (\ref{ex:text01-79})&\textit{Eori repaa kisi ni bono kehaa te ta naono.}&Then they tied the shell to the piece of wood.\\
    (\ref{ex:text01-80})&\textit{Nomana a kuhi na von riori tea sitoa.}&Today they buy the grater in the store.\\
    (\ref{ex:text01-81})&\textit{Na paku maa riori bona aeana.}&It is made from iron.\\
    (\ref{ex:text01-82})&\textit{A kuhi na kuhi varutarutaa riori bona overe.}&The grater is used to grate the coconut into small pieces.\\
    (\ref{ex:text01-83})&\textit{Meha maa vuri eori repaa vasibao me bari bona kaukau.}&Sometimes they \textit{vasibao} a sweet potato with it.\\
    (\ref{ex:text01-84})&\textit{Tea vasibao, ean na kuhi bata nom e ta overe bona kehaa.}&\textit{Vasibao} means that you grate a piece of coconut with a seashell.\\
    (\ref{ex:text01-85})&\textit{A kehaa vaa tea boon.}&A seashsell from the mangroves.\\
    (\ref{ex:text01-86})&\textit{A kehaa bona na dao riori bona boen.}&This seashell is called \textit{boen}.\\
    (\ref{ex:text01-87})&\textit{Ean pasi kuhi bata nom, bean ani bata me nom e bene ta kaukau.}&You grate (the coconut flesh) while eating it with a piece of sweet potato.\\
\end{xltabular}
\vspace{-.8cm}
\begin{figure}[H]
    \captionsetup[subfigure]{oneside,margin={0cm,0cm}}
    \captionsetup{justification=centering}
    \subfloat[]{\includegraphics[trim=0 20 0 85,clip,width=.43\linewidth]{figures/5_28.jpg}\label{fig:5_28}}\hfil
    \subfloat[]{\includegraphics[trim=0 0 0 0,clip,width=.57\linewidth]{figures/5_29.jpg}\label{fig:5_29}}
    \caption{(a) \textit{Varomo tope} -- Stacking up coconut shells,\\ (b) \textit{Vararana tope} -- Rows of coconut shells}
\end{figure}
% \begin{figure}[H]
%     \includegraphics[width=.5\linewidth]{figures/5_28.jpg}
%     \caption{\textit{Varomo tope} -- Stacking up coconut shells}
%     \label{fig:5_28}
% \end{figure}
% \begin{figure}[H]
%     \includegraphics[width=.5\linewidth]{figures/5_29.jpg}
%     \caption{\textit{Vararana tope} -- Rows of coconut shells}
%     \label{fig:5_29}
% \end{figure}
\begin{xltabular}{\textwidth}{lXX}
    (\ref{ex:text01-88})&\textit{\textbf{O tope}}&\textbf{The coconut shells}\\
    (\ref{ex:text01-89})&\textit{O tope na vaanoto suraa ni riori tea vamahaka kavuraa.}&The coconut shells are used for lighting the fire to dry copra.\\
    (\ref{ex:text01-90})&\textit{Varomo ni rori bono tope,}&They stack up the coconut shells,\\
    (\ref{ex:text01-91})&\textit{eori repaa vararana ni nao bari komana inu vaamahaka kavuraa.}&then they arrange them in rows in the copra drying house.\\
    (\ref{ex:text01-92})&\textit{O buaku geo kukan o vararana tope.}&There are two or three rows of coconut shells.\\
    (\ref{ex:text01-93})&\textit{Eori repaa vaanoto bono suraa matano vararana tope.}&Then they light the fire at the front of the row of coconut shells.\\
    (\ref{ex:text01-94})&\textit{O tope repaa asi petepete komana nao,}&The coconut shells then catch fire one after the other,\\
    (\ref{ex:text01-95})&\textit{ore suguna be asi vakavara.}&until they all burn.\\
    (\ref{ex:text01-96})&\textit{Meha maa vuri, eori repaa vaanoto bono suraa tope,}&Sometimes they light a fire of coconut shells,\\
    (\ref{ex:text01-97})&\textit{Be asi vakavara ge asi kavu,}&and when it has burnt down or burnt to ashes,\\
    (\ref{ex:text01-98})&\textit{eori repaa pupuhu iana toon no orana tope.}&they roast fish on the charcoals of the coconut shells.\\
\end{xltabular}
\begin{figure}[H]
    \includegraphics[trim=0 50 0 20,clip,width=.35\linewidth]{figures/5_30.jpg}\hfil
    \includegraphics[width=.65\linewidth]{figures/5_31.jpg}
    \caption{\textit{Pihae} -- Coconut frond}
    \label{fig:5_30-31}
\end{figure}
\begin{xltabular}{\textwidth}{lXX}
    (\ref{ex:text01-99})&\textit{\textbf{O pihae}}&\textbf{The coconut frond}\\
    (\ref{ex:text01-100})&\textit{O pihae vai o pereperee na paku riori bona bero a taba.}&The green coconut frond is used for many things.\\
    (\ref{ex:text01-101})&\textit{(1) Vin rori bona bonaa, mene dee tahii,}&(1) They weave the fishing basket for taking it to the sea,\\
    (\ref{ex:text01-102})&\textit{beori nao tea vagana tea masi tea hovuhovu iobo.}&when they go fishing on the reef to dig out sea-sausages.\footnote{See Chapter \ref{chap:39}.}\\
    (\ref{ex:text01-103})&\textit{(2) Vin rori bona hoi, mene vaobete tabaan vai a mahaka.}&(2) They weave baskets for storing cooked food.\\
    (\ref{ex:text01-104})&\textit{Amaa meha maa hoi vateen tagihu riori mohina,}&Other baskets they carry in backpacks to the garden in the jungle,\\
    (\ref{ex:text01-105})&\textit{eori repaa vaon maa bari bona tabaan vai a peree}&then they load them with the raw food\\
    (\ref{ex:text01-106})&\textit{eori re paa vateen voosu maa bari.}&and carry them home in backpacks.\\
\end{xltabular}
\begin{figure}[H]
    \captionsetup[subfigure]{oneside,margin={0cm,0cm}}
    \subfloat[]{\includegraphics[trim=0 0 0 20,clip,width=.45\linewidth]{figures/5_33.jpg}\label{fig:5_33}}
    \subfloat[]{\includegraphics[trim=0 0 0 0,clip,width=.55\linewidth]{figures/5_34.jpg}\label{fig:5_34}}
    \caption{(a) \textit{Vin hoi} -- Plaiting baskets, (b) \textit{Hoi} -- Basket}
\end{figure}
\begin{figure}[H]
    \includegraphics[trim=0 240 80 70,clip,width=.5\linewidth]{figures/5_32.jpg}
    \caption{\textit{O bonaa hoohoo} -- A fishing basket woven from midribs of coconut leaves}
    \label{fig:5_32}
\end{figure}
\begin{figure}[H]
    \includegraphics[width=.7\linewidth]{figures/5_35.jpg}
    \caption{\textit{Vateen voosu} -- Carry sth. home in a backpack}
    \label{fig:5_35}
\end{figure}
\begin{xltabular}{\textwidth}{lXX}
    (\ref{ex:text01-107})&\textit{(3) Vin rori bona petata, mene vaobete bona menaga.}&(3) They weave the \textit{petata}-trays for putting the tapioca dumplings on it.\\
    (\ref{ex:text01-108})&\textit{Bero amaa taba to vin riori bono pihae vai o pereperee.}&Many things are woven from the green coconut frond.\\
    (\ref{ex:text01-109})&\textit{O sinoo no pihae vai o perepereee na gono riori,}&They get the green stalk of the coconut frond,\footnotemark\\
    (\ref{ex:text01-110})&\textit{eori repaa kahu bono sikiri mene paku bono veveo.}&and then they scrape the midribs of the leaves to make the broom.\\
    (\ref{ex:text01-111})&\textit{O veveo na veveo riori bona maa nou vaa burihina ae komana maa inu.}&With the broom they sweep the rubbish in the backyard and inside the houses.\\
\end{xltabular}\footnotetext{lit. `bone of the coconut frond'}
\begin{figure}[H]
    \captionsetup[subfigure]{oneside,margin={0cm,0cm}}
    \subfloat[]{\includegraphics[trim=0 30 0 0,clip,width=.44\linewidth]{figures/5_36.jpg}\label{fig:5_36}}
    \subfloat[]{\includegraphics[trim=0 0 0 0,clip,width=.56\linewidth]{figures/5_37.png}\label{fig:5_37}}
    \caption{(a) \textit{Petata} -- Tray for dumplings, (b) \textit{Veveo} -- Broom}
\end{figure}
\begin{figure}[H]
    \captionsetup[subfigure]{oneside,margin={0cm,0cm}}
    \subfloat[]{\includegraphics[trim=0 30 0 0,clip,width=.37\linewidth]{figures/5_38.jpg}\label{fig:5_38}}
    \subfloat[]{\includegraphics[trim=0 0 0 0,clip,width=.63\linewidth]{figures/5_39.jpg}\label{fig:5_39}}
    \caption{(a) \textit{Rom} -- Stem, (b) \textit{Pagara} -- Footstalk}
\end{figure}
\begin{xltabular}{\textwidth}{lXX}
    (\ref{ex:text01-112})&\textit{Beo pihae mahaka teo overe,}&When the coconut frond has dried up on the coconut palm,\\
    (\ref{ex:text01-113})&\textit{o pagara repaa takapa kahi bono rom,}&the footstalk comes off the stem,\\
    (\ref{ex:text01-114})&\textit{o pihae repaa kuu vo teo kasuana.}&and the coconut frond falls down to the ground.\\
    (\ref{ex:text01-115})&\textit{O pihae vai o mahaka na rava riori bene guu,}&The dry coconut frond is used to singe the pig,\\
    (\ref{ex:text01-116})&\textit{beori mei rori bari tea kaku.}&before they butcher it.\\
    (\ref{ex:text01-117})&\textit{O pagara no overe, ae o nanisi na vaanoto suraa pete ni riori.}&The footstalk of the coconut palm and the panicle are also used for lighting the fire.\\
\end{xltabular}
\begin{figure}[H]
    \includegraphics[trim=0 80 0 250,clip,width=.5\linewidth]{figures/5_40.jpg}\hfil
    \includegraphics[trim=0 80 0 250,clip,width=.5\linewidth]{figures/5_41.jpg}
    \caption{\textit{Rava guu} -- Singeing pigs}
    \label{fig:5_40-41}
\end{figure}
\vspace{-.1cm}
\begin{xltabular}{\textwidth}{lXX}
    (\ref{ex:text01-118})&\textit{\textbf{O pesupaava}}&\textbf{The fibrous husk}\\
    (\ref{ex:text01-119})&\textit{O pesupaava mee na vaanoto suraa ni riori.}&The fibrous husk is also used for lighting fires.\\
    (\ref{ex:text01-120})&\textit{Meha maa vuri, eori repaa paku bari bona kanono.}&Sometimes they make a string from it.\\
    (\ref{ex:text01-121})&\textit{Sirisiri varutarutaa riori,}&They tear it into small pieces,\\
    (\ref{ex:text01-122})&\textit{eori repaa taneo tea vinugu bari.}&then they start twisting a string from them (the fibres of the husk).\\
\end{xltabular}
\vspace{-.4cm}
\begin{figure}[H]
    \captionsetup[subfigure]{oneside,margin={0cm,0cm}}
    \captionsetup{justification=centering}
    \subfloat[]{\includegraphics[trim=0 0 0 30,clip,width=.56\linewidth]{figures/5_42.jpg}\label{fig:5_42}}
    \subfloat[]{\includegraphics[trim=0 0 0 0,clip,width=.44\linewidth]{figures/5_43.jpg}\label{fig:5_43}}
    \caption{(a) \textit{O pesupaava} -- The fibrous coconut husk,\\ (b) \textit{Sirisiri varutarutaa} -- Tearing into small pieces}
\end{figure}
\begin{xltabular}{\textwidth}{lXX}
    (\ref{ex:text01-123})&\textit{\textbf{O vaanomo}}&\textbf{The coconut oil}\\
    (\ref{ex:text01-124})&\textit{Beara tau paku o vaanomo, eara na nao rara tea hovu overe.}&When we are about to make oil, we go and husk coconuts.\\
    (\ref{ex:text01-125})&\textit{Beara rake o vaanomo vai o beera, eara repaa hovu amaa buaku saavun ge kukan saavun overe.}&When we want a lot of oil, we husk twenty or thirty coconuts.\\
    (\ref{ex:text01-126})&\textit{Beara hovu vakavara amaa overe, eara repaa taneo tea kuhi overe.}&When we have finished husking the coconuts, we start grating coconuts.\\
    (\ref{ex:text01-127})&\textit{Murinae eara repaaa posi overe vo tea nahu.}&After that we squeeze the (grated) coconut (so that the liquid flows) into a pot.\\
    (\ref{ex:text01-128})&\textit{Beara posi vakavara a overe, eara repaa koma hopo nio sii hum ruene tea nahu.}&When we have finished squeezing the coconut, we pour a little water into the pot.\\
    (\ref{ex:text01-129})&\textit{Eara repaa vahio a nahu teo suraa.}&Then we set the pot onto the fire.\\
    (\ref{ex:text01-130})&\textit{Be gonono, eara paa gono a sipuni, eara paa verete bata rara a overe.}&When it is boiling, we get a spoon and then keep stirring the coconut (liquid).\\
    (\ref{ex:text01-131})&\textit{Beara sa verete bata haara a overe, o konoboo repaa asi.}&If we did not keep stirring the coconut (liquid), the coconut milk would burn.\\
    (\ref{ex:text01-132})&\textit{Beara tara bea rikisi tavusu, o vaanomo repaa agaa vaavaha bata.}&When we see that the cream comes out, then the oil floats separately.\\
    (\ref{ex:text01-133})&\textit{Ean repaa vaarutaa bata o suraa.}&And you make the fire smaller.\\
    (\ref{ex:text01-134})&\textit{Bean tara beo vaanomo tean maraa.}&When you see that your oil is clear,\\
    (\ref{ex:text01-135})&\textit{ean repaa kae kahi e bono suraa.}&you lift it (the pot) from the fire.\\
    (\ref{ex:text01-136})&\textit{Beo vaanomo kunio, ean repaa isuvu e tea botori tean.}&When the oil is cold, you fill it into your bottle.\\
    (\ref{ex:text01-137})&\textit{Ean paa vaauhu nie,}&You rub it onto your skin,\\
    (\ref{ex:text01-138})&\textit{a suinman re vakasi oha.}&so that your skin shines all over.\\
    (\ref{ex:text01-139})&\textit{Bean rake o vaanomo tean re avuhu mataa,}&If you want your oil to smell good,\\
    (\ref{ex:text01-140})&\textit{ean paa osi maa o purupuru no naono to dao riori bono binaa.}&you pick the flowers of a tree that is called ylang-ylang.\\
    (\ref{ex:text01-141})&\textit{Ean paa varoava e.}&Then you put them in the sun.\\
    (\ref{ex:text01-142})&\textit{Be mahaka, ore paa avuavuhu.}&When they are dry, they smell good.\\
    (\ref{ex:text01-143})&\textit{Bean nahu o vaanomo tean, ean paa vaobete hopo nao o purupuru no binaa teo vaanomo, be gonono nana.}&When you cook your oil, you put the flowers of the ylang-ylang into the oil, when it is boiling.\\
    (\ref{ex:text01-144})&\textit{Meha papana na vaatei rori bona paka sukuita.}&Other people put \textit{sukuita} leaves (into the oil).\\
    (\ref{ex:text01-145})&\textit{Tei nana tea peepeha aba}&It is up to everybody\\
    (\ref{ex:text01-146})&\textit{tea viana bona paka naono tabae toro vaobete eve teo vaanomo,}&to select whatever leaves he or she is going to put into the oil,\\
    (\ref{ex:text01-147})&\textit{o vaanomo re avuhu mataa.}&so that the oil smells good.\\
    (\ref{ex:text01-148})&\textit{Eori he toro vaobete vai koa bari teo vuri nahu.}&But they must only put them into the oil when they cook them.\\
\end{xltabular}
\begin{figure}[H]
    \captionsetup[subfigure]{oneside,margin={0cm,0cm}}
    \captionsetup{justification=centering}
    \subfloat[]{\includegraphics[trim=0 0 0 0,clip,width=.49\linewidth]{figures/5_44.jpg}\label{fig:5_44}}
    \subfloat[]{\includegraphics[trim=15 0 0 0,clip,width=.51\linewidth]{figures/5_45.jpg}\label{fig:5_45}}
    \caption{(a) \textit{A inu vaamahaka kavuraa} -- Copra drying house,\\ (b) \textit{Vararana tope} -- Rows of coconut shells}
\end{figure}
\begin{xltabular}{\textwidth}{lXX}
    (\ref{ex:text01-149})&\textit{\textbf{Kavuraa}}&\textbf{Copra}\\
    (\ref{ex:text01-150})&\textit{Beara tau paku a kavuraa,}&When we are about to make copra,\\
    (\ref{ex:text01-151})&\textit{taba momohu toro paku raara,}&the first thing we must do,\\
    (\ref{ex:text01-152})&\textit{eara toro rapisi roho a garasi  vaa komana overe.}&we must cut the grass in the coconut plantation.\\
    (\ref{ex:text01-153})&\textit{Eara repaa taneo tea unuunu overe.}&Then we start collecting coconuts.\\
    (\ref{ex:text01-154})&\textit{Bean unuunu vakavara a overe, ean repaa taneo tea hovu.}&When you have finished collecting the coconuts, you start husking them.\footnote{Note that \textit{a overe} lit. `the coconut’ in (154), (156), (158), (160) is the singular form, but translated by the English plural form.}\\
    (\ref{ex:text01-155})&\textit{Bean hovu vakavara, ean re paa taneo tea booboha.}&When you have finished husking (them), you start breaking (them).\\
    (\ref{ex:text01-156})&\textit{Bean booboha vakavara, ean repaa deedee a overe to paa booboha vuan vo inu vaamahaka kavuraa.}&When you have finished breaking (them), you carry the coconuts that you have broken to the copra drying house.\\
\end{xltabular}
\begin{figure}[H]
    \includegraphics[trim=0 50 0 100,clip,width=.8\linewidth]{figures/5_46.jpg}
    \caption{\textit{Tahi kavuraa} -- Pushing the copra into bags with a stick}
    \label{fig:5_46}
\end{figure}
\begin{xltabular}{\textwidth}{lXX}
    (\ref{ex:text01-157})&\textit{Be kavara a deedee overe, ean repaa tasu hiava nie tea sibono.}&When all the coconuts have been carried there, you throw them up onto the roast.\\
    (\ref{ex:text01-158})&\textit{Bean tasu hiava vakavara nia overe tea sibono,}&When you have finished throwing the coconuts onto the roast,\\
    (\ref{ex:text01-159})&\textit{ean repaa taneo tea vaakopo.}&you start arranging (them) upside down.\\
    (\ref{ex:text01-160})&\textit{Bean vaakopo vakavara a overe tea sibono,}&When you have finished arranging the coconuts on the roast,\\
    (\ref{ex:text01-161})&\textit{ean repaa taneo tea vararana tope paana sibono kavuraa.}&you start placing coconut shells in rows underneath the copra roast.\\
    (\ref{ex:text01-162})&\textit{O kukan o vararana to paku raara paana sibono komana inu sibono kavuraa.}&Three rows we make underneath the roast in the copra drying house.\\
    (\ref{ex:text01-163})&\textit{Ean repaa vaanoto suraa matano vararana tope.}&Then you light the fire at the front of the row of coconut shells.\\
    (\ref{ex:text01-164})&\textit{Pasi asiu ore suguna beo vararana tope asi vakavara.}&It will catch fire until the whole row of coconut shells burns.\\
    (\ref{ex:text01-165})&\textit{O kukan geo dorana o boni to tei nae tea sibono,}&Three or four days they stay on the roast,\\
    (\ref{ex:text01-166})&\textit{bean vaanoto nom.}&when you light (the fire).\\
    (\ref{ex:text01-167})&\textit{Bean tara be ta kavuraa takapa kahi bona tope,}&When you see that a piece of copra is coming off the shell,\\
    (\ref{ex:text01-168})&\textit{ean paa kona voen, a kavuraa tean a mahaka.}&then you know that your copra is dry.\\
    (\ref{ex:text01-169})&\textit{Bea kavuraa tean mahaka,}&When your copra is dry,\\
    (\ref{ex:text01-170})&\textit{ean repaa kapa kahi e bono tope.}&you remove it from the shells.\\
    (\ref{ex:text01-171})&\textit{Bean kapa vakavara a kavuraa tean,}&When you have removed your copra,\\
    (\ref{ex:text01-172})&\textit{ean repaa mosi varutarutaa e.}&you cut it into small pieces.\\
    (\ref{ex:text01-173})&\textit{Murinae ean repaa baeke eve, bean tahi bata nom e.}&After that you put them into bags by pushing them with a stick.\\
    (\ref{ex:text01-174})&\textit{Bean tahi vahapusu o baeke kavuraa tean,}&When you have got your copra bag  fully stuffed,\\
    (\ref{ex:text01-175})&\textit{ean repaa kuna pina nia matano baeke kavuraa.}&you sew up the opening of the copra bag.\\
    (\ref{ex:text01-176})&\textit{A vagegeana kiu, bean vaahae o baeke kavuraa tean tea kaara,}&The last work is that you load your bag of copra on a truck,\\
    (\ref{ex:text01-177})&\textit{ean paa nao tea vaavon Buka.}&and then you go to sell it in Buka.\\
\end{xltabular}
\begin{xltabular}{\textwidth}{lXX}
    (\ref{ex:text01-178})&\textit{\textbf{O rom overe}}&\textbf{Coconut palm stem}\\
    (\ref{ex:text01-179})&\textit{O rom overe na romrom vakakato riori,}&They cut the stem of the coconut palm crosswise into short logs,\\
    (\ref{ex:text01-180})&\textit{eori repaa kepo bari mene vaanoto suraa vaa tea taotao ae tea nahunahu kaukau gea tabae to rake riori tea nahu.}&then they chop it in order to light the fire to bake (food) in the earth oven and cook sweet potatoes in pots or whatever they want to cook.\\
    (\ref{ex:text01-181})&\textit{Meha maa vuri o rom overe na kopaa riori,}&Sometimes they cut the stem of a coconut palm lengthwise,\\
    (\ref{ex:text01-182})&\textit{eori repaa paku bari bono inaho.}&to make the taro planting stick from it.\\
    (\ref{ex:text01-183})&\textit{Kopaa riori, eori repaa taraha tatagi kurusu bari.}&They cut it lengthwise, and then they very carefully carve the surface.\\
    (\ref{ex:text01-184})&\textit{O inaho toro raramun,}&The planting stick must be smooth,\\
    (\ref{ex:text01-185})&\textit{ahiki ta taba re tou bona kuri raara, ae ruvis bona.}&so that there is not anything (that) would prick our hand and scratch it.\\
\end{xltabular}
\begin{figure}[H]
    \includegraphics[width=.55\linewidth]{figures/5_47.jpg}
    \caption{\textit{O inaho} -- The planting stick}
    \label{fig:5_47}
\end{figure}
\end{document}